% DRAIN OUT - Pitch Deck Draft: the content of every slide
% Compile on Overleaf with pdfLaTeX. No external files needed.

\documentclass[11pt,a4paper]{article}

\usepackage[margin=1.9cm]{geometry}
\usepackage[T1]{fontenc}
\usepackage{lmodern}
\usepackage{microtype}
\usepackage{amsmath,amssymb}
\usepackage{booktabs,tabularx,array,longtable}
\usepackage{enumitem}
\usepackage[table]{xcolor}
\usepackage[most]{tcolorbox}
\usepackage{titlesec}
\usepackage{fancyhdr}
\usepackage[colorlinks=true,linkcolor=accent,urlcolor=accent]{hyperref}

% ---------- Look ----------
\definecolor{accent}{HTML}{0B5CAD}
\definecolor{hook}{HTML}{C0392B}
\definecolor{visual}{HTML}{1E7B34}
\definecolor{speak}{HTML}{6C3483}
\definecolor{soft}{HTML}{EEF3FA}
\definecolor{ink}{HTML}{1B2330}

\titleformat{\section}{\Large\bfseries\color{accent}}{\thesection}{0.8em}{}
\titleformat{\subsection}{\large\bfseries\color{ink}}{\thesubsection}{0.7em}{}

\setlength{\headheight}{14pt}
\pagestyle{fancy}
\fancyhf{}
\lhead{\small\color{accent}\textbf{DRAIN OUT} \textperiodcentered{} Pitch Deck Draft}
\rhead{\small Team Tech Sprints \textperiodcentered{} SIH 2026}
\cfoot{\small\thepage}
\renewcommand{\headrulewidth}{0.4pt}

\setlist{itemsep=2pt,topsep=3pt}
\renewcommand{\arraystretch}{1.3}
\newcolumntype{L}{>{\raggedright\arraybackslash}X}

% One box per slide.
\newtcolorbox{slidebox}[1]{enhanced,breakable,colback=white,colframe=accent,
  coltitle=white,fonttitle=\bfseries\large,title={#1},arc=2pt,boxrule=0.9pt,
  left=9pt,right=9pt,top=7pt,bottom=7pt,before skip=12pt,after skip=12pt}

% The parts of a slide.
\newcommand{\purpose}[1]{\par{\small\color{ink!70}\textbf{PURPOSE:} \emph{#1}}\par\medskip}
\newcommand{\hookline}[1]{%
  \begin{tcolorbox}[colback=hook!6,colframe=hook,boxrule=0.6pt,arc=2pt,left=7pt,right=7pt,top=4pt,bottom=4pt]
  {\color{hook}\textbf{HEADLINE (top of slide):}} \textbf{#1}\end{tcolorbox}}
\newcommand{\slidepart}[2]{\par\smallskip{\color{#1}\textbf{#2}}\par\smallskip}
\newcommand{\onslide}{\slidepart{accent}{ON THE SLIDE}}
\newcommand{\visuals}{\slidepart{visual}{VISUALS \& LAYOUT}}
\newcommand{\speaker}{\slidepart{speak}{SPEAKER SCRIPT}}
\newcommand{\say}[1]{\begin{tcolorbox}[colback=speak!5,colframe=speak!60,boxrule=0.4pt,arc=2pt,
  left=7pt,right=7pt,top=4pt,bottom=4pt]\itshape #1\end{tcolorbox}}

% =====================================================================
\begin{document}

\begin{titlepage}
  \centering
  \vspace*{3cm}
  {\color{accent}\rule{\linewidth}{1.2pt}}\par\vspace{0.6cm}
  {\Huge\bfseries DRAIN OUT\par}
  \vspace{0.3cm}
  {\LARGE Pitch Deck Draft\par}
  \vspace{0.4cm}
  {\large The content, visuals and script of every slide\par}
  \vspace{0.6cm}
  {\color{accent}\rule{\linewidth}{1.2pt}}\par
  \vspace{1.2cm}
  {\Large\itshape ``Know the flood before it reaches your street.''\par}
  \vspace{1.5cm}
  \begin{tabular}{@{}ll@{}}
    \textbf{Event}     & Smart India Hackathon 2026 \\
    \textbf{Team}      & Tech Sprints \\
    \textbf{Prototype} & \url{https://github.com/QuantaBoy/SIH26085} \\
  \end{tabular}
  \vfill
  \begin{minipage}{0.88\linewidth}\small
  \textbf{How to read this.} Each slide has five parts:
  {\color{ink!70}\textbf{Purpose}} (what the slide must make the jury believe),
  {\color{hook}\textbf{Headline}} (the one line at the top),
  {\color{accent}\textbf{On the slide}} (the exact text),
  {\color{visual}\textbf{Visuals \& layout}} (what to draw and where), and
  {\color{speak}\textbf{Speaker script}} (what to say while it is on screen).
  \end{minipage}
\end{titlepage}

\tableofcontents
\newpage

% =====================================================================
\section{The story the deck tells}

\begin{center}
\begin{tabular}{@{}clp{9.5cm}@{}}
\toprule
\textbf{Slide} & \textbf{Beat} & \textbf{What the jury should believe after it}\\
\midrule
1 & Identity & Who we are and which problem we solve\\
2 & Problem $\rightarrow$ Solution & Warnings stop at the ward; DRAIN OUT goes to the street\\
3 & How it works & A rigorous physics + AI pipeline, not a dashboard on top of a weather API\\
4 & Can it be done & Every input exists nationally; every risk has a mitigation\\
5 & Why it matters & Lives, money and city operations, for named users\\
6 & Evidence & We know the field, and we fill the gap it leaves\\
\bottomrule
\end{tabular}
\end{center}

\noindent\textbf{Tagline (use on slides 1 and 5, and as the last words of the pitch):}
\emph{Know the flood before it reaches your street.}

% =====================================================================
\section{Slide by slide}

% ---------------------------------------------------------------------
\begin{slidebox}{Slide 1 \textendash{} Title Page}
\purpose{Identify the team and problem instantly; set a confident, national tone.}
\hookline{DRAIN OUT: Street-Level Urban Flood Nowcasting \& Safe Routing for Indian Cities}

\onslide
\begin{itemize}
  \item \textbf{Problem Statement ID:} \emph{(fill in)}
  \item \textbf{Problem Statement Title:} \emph{(fill in)}
  \item \textbf{Theme:} \emph{(fill in)}
  \item \textbf{PS Category:} Software
  \item \textbf{Team ID:} \emph{(fill in)}
  \item \textbf{Team Name:} Tech Sprints
\end{itemize}
Tagline under the title: \emph{Know the flood before it reaches your street.}

\visuals
\begin{itemize}
  \item Keep the SIH template; put the DRAIN OUT logo beside the title.
  \item No field may be left empty: an empty title page is the first thing a jury notices.
\end{itemize}

\speaker
\say{``We are Tech Sprints. Our solution is DRAIN OUT: it tells you which street will flood,
how deep, and when, and gets you around it.''}
\end{slidebox}

% ---------------------------------------------------------------------
\begin{slidebox}{Slide 2 \textendash{} Proposed Solution \& Unique Value Proposition}
\purpose{Make the problem felt in one sentence, then show a solution that is clearly
different from what exists.}
\hookline{Today's flood warnings say ``Chennai''. DRAIN OUT says ``Anna Salai, 35 cm, 4:40 PM''.}

\onslide
\textbf{Proposed Solution.} DRAIN OUT is a national platform that forecasts urban
flooding street by street, 0--3 hours ahead, and guides people and emergency services
around it.
\begin{itemize}
  \item \textbf{Coupled drainage + terrain simulation} predicts flood depth in
        \textbf{cm on every street, every 5 minutes}.
  \item \textbf{Onboards any Indian city}: drainage maps (GIS, CAD, scanned) become a
        working drain network, with missing links auto-repaired.
  \item \textbf{Physics-trained AI (ConvLSTM)} with a \textbf{historical pattern model}:
        millisecond forecasts that learn each city's recurring flood zones.
  \item Detects \textbf{overflowing manholes and overloaded drains}, turning fixed
        desilting schedules into targeted ones.
  \item \textbf{Flood-safe routing (A* / Dijkstra)} for citizens, ambulances and traffic police.
  \item Delivered through a \textbf{control-room GIS dashboard, a multilingual mobile app
        and open APIs}, integrated with NDMA alerts.
\end{itemize}

\textbf{Unique Value Proposition} (six boxes around the logo):
\begin{tabularx}{\linewidth}{@{}LL@{}}
1.\ Street-level, not ward-level: depth in cm & 4.\ Manhole-level overflow and backflow detection\\
2.\ Missing-link-aware drainage network & 5.\ One platform, any Indian city\\
3.\ Physics + AI + historical memory & 6.\ Flood-safe routing in under a second\\
\end{tabularx}

\textbf{DRAIN OUT addresses the problem by} (four icon boxes):
\begin{itemize}
  \item Street-level flood depth (cm), 0--3 hours ahead, every 5 minutes
  \item Pinpointing overloaded drains and overflowing manholes before they overflow
  \item Rerouting traffic and emergency vehicles around flooded roads
  \item Location-based early warnings in local languages
\end{itemize}

\visuals
\begin{itemize}
  \item Left half: Proposed Solution. Right half: the six UVP boxes around the DRAIN OUT logo.
  \item Bottom strip: \textbf{Before / After} \textendash{} ``Flood-unaware route'' (red line
        through a flooded street) $\rightarrow$ ``Flood-aware safe route'' (green line around it).
  \item Show the dashboard on a laptop and a phone; label the phone ``Mobile app''.
\end{itemize}

\speaker
\say{``It's 4 PM. The alert on your phone says: heavy rain expected in Chennai. An ambulance
driver reads it too. Which road is safe? Which underpass floods first? The alert can't
say. Every flood warning in India today stops at the city or ward level, but water does
not flood a ward, it floods a street. DRAIN OUT forecasts flooding street by street, in
centimetres, three hours ahead, every five minutes, and routes you around it.''}
\end{slidebox}

% ---------------------------------------------------------------------
\begin{slidebox}{Slide 3 \textendash{} Technical Approach}
\purpose{Prove technical depth: real engineering and AI, in a clean pipeline the jury can
follow left to right.}
\hookline{Rain, drains and roads \textendash{} simulated together, every 5 minutes.}

\onslide
\begin{enumerate}[label=\textbf{\arabic*.}]
  \item \textbf{Data collection:} IMD Doppler radar and rain gauges \textperiodcentered{}
        CartoDEM / LiDAR \textperiodcentered{} city drainage records \textperiodcentered{}
        NRSC land use and soil (SCS-CN) \textperiodcentered{} INCOIS tide levels
        \textperiodcentered{} historical floods (NDMA, Sentinel-1) \textperiodcentered{}
        IoT water-level sensors
  \item \textbf{City onboarding:} digitise $\rightarrow$ build drain network $\rightarrow$
        auto-repair gaps $\rightarrow$ infer flow direction $\rightarrow$ correct terrain
        to surveyed levels
  \item \textbf{Physics layer:} SCS-CN runoff $\rightarrow$ drain capacity
        $Q_{cap} = \frac{1}{n} A R^{2/3} S^{1/2}$ $\rightarrow$ drains and street surface
        simulated together $\rightarrow$ overflow and backflow at manholes
  \item \textbf{AI layer:} ConvLSTM trained on simulated storms and historical floods
        \textperiodcentered{} historical pattern model \textperiodcentered{} live correction
        from sensors \textperiodcentered{} flood probability per street
  \item \textbf{Decision support:} early warnings \textperiodcentered{} desilting priority
        \textperiodcentered{} diversions \textperiodcentered{} flood-safe routing
  \item \textbf{Delivery:} GIS dashboard (ICCC) \textperiodcentered{} mobile app
        \textperiodcentered{} REST APIs \textperiodcentered{} NDMA SACHET (CAP) alerts
\end{enumerate}
\textbf{Outputs:} street-level depth (cm) \textperiodcentered{} 0--3 h nowcast at
5-minute steps \textperiodcentered{} flood risk and severity \textperiodcentered{}
overflowing manholes \textperiodcentered{} flood-safe routes

\textbf{Tech stack:} Python \textperiodcentered{} FastAPI \textperiodcentered{}
NumPy / SciPy \textperiodcentered{} PyTorch (ConvLSTM) \textperiodcentered{} PostGIS
\textperiodcentered{} Leaflet / GIS \textperiodcentered{} Docker

\textbf{Prototype:} working city-wide pilot on Chennai \textendash{}
\url{https://github.com/QuantaBoy/SIH26085} \textperiodcentered{} Demo video: \emph{(link)}

\visuals
\begin{itemize}
  \item A left-to-right flow diagram with six numbered stages, matching the list above.
  \item Inside the physics layer, show the Manning formula: it signals engineering rigour.
  \item Two yellow output boxes: \emph{Physics outputs} (depth, 2D surface flow, drain
        capacity, runoff, surcharge / backflow) and \emph{Final outputs} (street depth in
        cm, 0--3 h at 5 min, risk and severity).
  \item Right column: tech stack, GitHub link with the GitHub icon, demo video link.
  \item A small screenshot of the dashboard showing a real forecast with red streets
        \textendash{} never a screenshot with an error on it.
\end{itemize}

\speaker
\say{``Rainfall comes from IMD radar. The terrain and every drain in the city become one
network. Our physics engine simulates rain, drains and roads together, so it knows which
manhole overflows first. On top, a ConvLSTM trained on thousands of simulated storms
makes it fast enough to run every city, and a historical model remembers where floods
recur. Everything ends in one output: which street, how deep, and when.''}
\end{slidebox}

% ---------------------------------------------------------------------
\begin{slidebox}{Slide 4 \textendash{} Feasibility and Viability}
\purpose{Remove doubt: every input exists nationally, every risk has an answer, and it
scales without being rebuilt.}
\hookline{Any Indian city goes live from its own drainage records \textendash{} no new code.}

\onslide
\textbf{Feasibility Analysis}

\begin{tabularx}{\linewidth}{@{}p{3.2cm}L@{}}
\toprule
\textbf{Data} & Every input has a national source: IMD, NRSC, Copernicus, INCOIS, NDMA, city drainage records.\\
\textbf{Modelling} & Built on established engineering methods: Manning, SCS-CN, CPHEEO, HEC-22.\\
\textbf{AI} & Trained on simulated storms, so it does not depend on scarce flood records.\\
\textbf{Scalability} & The same pipeline for every city; a new city is data onboarding, not new code.\\
\textbf{Integration} & Open APIs connect to Smart City command centres, 108, traffic police and NDMA alerts.\\
\bottomrule
\end{tabularx}

\medskip
\textbf{Viability: Challenge $\rightarrow$ Mitigation}

\begin{tabularx}{\linewidth}{@{}p{4.4cm}L@{}}
\toprule
\textbf{Incomplete drainage records} & Auto-repair of missing links plus digitisation of paper maps\\
\textbf{Validation} & Each city tested against its past floods before go-live\\
\textbf{Model drift over years} & Recalibrated after every monsoon; historical model refitted on recent years\\
\textbf{Multi-agency integration} & Standard REST and CAP alert formats\\
\textbf{Nationwide scale} & Cloud deployment; one platform serving many cities\\
\bottomrule
\end{tabularx}

\medskip
\textbf{Viability statement:} DRAIN OUT is viable because it combines nationally
available data with established hydraulic engineering and AI in one modular pipeline
that is validated on each city's past floods and scaled city by city.

\visuals
\begin{itemize}
  \item Left: five blue feasibility labels with their text boxes.
  \item Right: five challenges, each joined by an orange icon to its mitigation.
  \item Bottom: the viability statement in a bordered box.
  \item Replace ``Your Team Name'' in the template's oval with ``Tech Sprints''.
\end{itemize}

\speaker
\say{``Every input already exists across India: IMD radar, NRSC land data, national
terrain, and every city's own drainage records. Where records are incomplete, and they
always are, our network repairs missing links automatically. The AI learns from
simulations, so it doesn't wait for decades of flood records. And a new city is a data
job, not a rebuild.''}
\end{slidebox}

% ---------------------------------------------------------------------
\begin{slidebox}{Slide 5 \textendash{} Impact and Benefits}
\purpose{Make the value concrete: who is helped, how, and why it matters nationally.}
\hookline{From a flood warning to a safe route \textendash{} in under a second.}

\onslide
\begin{tabularx}{\linewidth}{@{}p{3.2cm}L@{}}
\toprule
\textbf{01 Safety} & Street-level warnings up to 3 hours ahead; flood-safe routes for citizens, ambulances and fire services.\\
\textbf{02 Economic} & Less property and vehicle damage; desilting budgets spent where drains actually overflow.\\
\textbf{03 Urban planning} & Recurring flood hotspots and undersized drains identified for long-term investment.\\
\textbf{04 Operational} & One shared, live picture for city corporations, disaster authorities, traffic police and hospitals.\\
\bottomrule
\end{tabularx}

\medskip
\textbf{Core impact:} Street-level early warning and safe mobility for India's
flood-prone cities.

\textbf{Users:} city corporations \textperiodcentered{} NDMA / state disaster authorities
\textperiodcentered{} traffic police \textperiodcentered{} 108 ambulance and fire
services \textperiodcentered{} citizens

\textbf{SDGs:} 3 Good Health \textperiodcentered{} 9 Industry \& Infrastructure
\textperiodcentered{} 11 Sustainable Cities \textperiodcentered{} 13 Climate Action

\visuals
\begin{itemize}
  \item Keep the four-quadrant layout (Safety, Economic, Urban, Operational) around the
        SIH logo.
  \item Core impact in a bordered box at the bottom left; SDG icons at the bottom right.
  \item Add the user icons (municipal building, police, ambulance, citizen) in one row.
\end{itemize}

\speaker
\say{``For citizens, a warning that names their street and a way around the water. For
ambulances and fire services, routes that never drive into a flood. For city
corporations, the drains that will overflow, so they are cleaned before the monsoon, not
after. Know the flood before it reaches your street.''}
\end{slidebox}

% ---------------------------------------------------------------------
\begin{slidebox}{Slide 6 \textendash{} Research and References}
\purpose{Show we know the field and that DRAIN OUT fills a real gap in it.}
\hookline{Existing systems stop at the ward. We go to the street.}

\onslide
{\small
\begin{tabularx}{\linewidth}{@{}p{3.2cm}p{4.0cm}L@{}}
\toprule
\textbf{Organisation} & \textbf{Work} & \textbf{Gap DRAIN OUT fills}\\
\midrule
Ministry of Earth Sciences & I-FLOWS Mumbai & Ward-level alerts $\rightarrow$ street-level depth in cm\\
NCCR & C-FLOWS Chennai & Precomputed scenarios $\rightarrow$ live simulation every 5 min\\
NDMA & Urban Flooding Guidelines (2010) & Guidelines $\rightarrow$ an operational forecasting tool\\
MoHUA / CPHEEO & Storm Water Drainage Manual (2019) & Design standard $\rightarrow$ real-time drain load\\
FHWA & HEC-22 Urban Drainage Design & Inlet capture applied to live forecasting\\
IMD & Doppler Weather Radar network & Radar rainfall $\rightarrow$ street-level flood depth\\
\bottomrule
\end{tabularx}}

\visuals
\begin{itemize}
  \item Left: the table, with each organisation's logo in the first column and a working
        link for every work.
  \item Right: the two charts (urban flood deaths by city; area affected by floods) with
        a \textbf{source line under each}.
  \item Drop any paper that is not about urban flooding.
\end{itemize}

\speaker
\say{``India's flood warning systems are strong at the city and ward scale. I-FLOWS gives
ward-level alerts; C-FLOWS matches precomputed scenarios. DRAIN OUT fills the gap they
leave: live simulation, down to the street and the manhole, with a safe way around.''}
\end{slidebox}

% =====================================================================
\section{Jury questions and answers}

{\small
\begin{longtable}{@{}p{5.0cm}p{11.3cm}@{}}
\toprule
\textbf{Question} & \textbf{Answer}\\
\midrule\endhead
What have you built so far? & A working city-wide pilot on Chennai: the drain network, the
  coupled physics engine, the street-level forecast and flood-safe routing, running on real
  data. The ConvLSTM layer and IoT integration are the next phase.\\
Where does ConvLSTM get training data when floods are rare? & From the physics engine: it
  simulates thousands of storms, and the network learns to reproduce them in milliseconds.
  Scarce flood records are used to validate, not to train.\\
How accurate is it? & Every run conserves water exactly. Each city is validated against its
  past floods before go-live; for Chennai, the December 2015 flooding points.\\
Cities have incomplete drainage maps. & Our onboarding digitises paper maps and repairs
  missing links and unsized drains automatically, marking them as assumed.\\
Why not use I-FLOWS or C-FLOWS? & They warn at ward level or match precomputed scenarios.
  DRAIN OUT simulates live, street by street and manhole by manhole, and adds routing.\\
How does it scale to all of India? & One pipeline. A new city needs its drainage records,
  national terrain and IMD rainfall; no new code.\\
What about blocked drains? & Each drain's capacity is reduced by its surveyed condition, and
  a silting setting shows what happens if drains are partly blocked.\\
Who uses it and who pays? & City corporations and Smart City command centres; traffic police
  and 108 for routing. Built on free national data, so the cost is deployment, not licences.\\
How fast is it? & A city-wide 3-hour forecast in seconds; a flood-safe route in under a
  second.\\
\bottomrule
\end{longtable}}

% =====================================================================
\section{Final checklist before submission}

\begin{itemize}[label=$\square$]
  \item Title page: every field filled in.
  \item No template placeholder left (``Your Team Name'', ``LINK'').
  \item Prototype link and demo video link on slide 3 work.
  \item Dashboard screenshot shows a real forecast, no error messages.
  \item Every chart has a source line.
  \item Every reference is about urban flooding and has a working link.
  \item Spelling: ``Addresses'', ``maintenance'', ``permeability''.
  \item Each slide has its headline at the top.
  \item The pitch has been rehearsed to 3 minutes, ending on the tagline.
\end{itemize}

\end{document}
